% TOP_PROBLEM: {"id": 30006860, "title": "Chowla Non-Vanishing Conjecture for Dirichlet L-Functions at the Central Point", "category_id": 1, "category": {"id": 1, "name": "number_theory", "display_name": "Number Theory", "description": "Properties of integers, prime numbers, Diophantine equations.", "slug": "number-theory", "order_index": 1, "created_at": "2026-07-31T15:26:25.670Z"}, "status": "open"}
% FIRST_POSTED: 2026-09-27
% !TeX program = lualatex
% ATTEMPT_STATUS: unresolved
\documentclass[11pt]{article}
\usepackage[margin=25mm]{geometry}
\usepackage{fontspec}
\setmainfont{Latin Modern Roman}
\usepackage{amsmath,amssymb,mathtools}
\usepackage{unicode-math}
\setmathfont{Latin Modern Math}
\usepackage{xurl,hyperref}
\hypersetup{colorlinks=true,urlcolor=blue}
\setcounter{secnumdepth}{0}
\setlength{\parindent}{0pt}
\setlength{\parskip}{5pt}
\setlength{\emergencystretch}{3em}
\begin{document}
\section*{\texorpdfstring{Chowla Non-Vanishing Conjecture for Dirichlet L-Functions at the Central Point}{Chowla Non-Vanishing Conjecture for Dirichlet L-Functions at the Central Point}}\label{30006860}
\subsection{Definitions and mathematical statement}
A Dirichlet character modulo $q$ is a homomorphism $({\mathbb{Z}}/q{\mathbb{Z}})^\times\to{\mathbb{C}}^\times$, extended periodically to integers and by zero on nonunits. It is primitive of conductor $q$ if it is not induced from a proper divisor of $q$. Initially, for $\Re s>1$,
\[
 L(s,\chi)=\sum_{n\ge1}\frac{\chi(n)}{n^s};
\]
use its analytic continuation to define the central value. Chowla's nonvanishing assertion in the selected scope is $L(\tfrac12,\chi)\ne0$ for every primitive character, real or complex, of either parity $\chi(-1)=\pm1$. The conductor-one case is the zeta function. No positivity or uniform lower bound is demanded. Primitivity matters: the statement does not simply extend the quantifier to arbitrary imprimitive characters with additional Euler factors.
\par\smallskip\textit{Formulation and concept references:} \href{https://arxiv.org/pdf/2603.22124v1}{[S1]}.

\subsection{Short English statement}
Is the central value of every primitive Dirichlet L-function nonzero, regardless of its character's parity or whether the character is real?
\subsection{Research attempt: signed blocks, root numbers, and finite certificates}
\paragraph{1. Scope and source audit.}
The quantifier is every primitive Dirichlet character, including real and complex characters of either parity and the conductor-one character. Earnst's March 2026 paper [S1], read in its HTML version [E1] on 27 September 2026, proves positive-proportion nonvanishing for even primitive characters of prime conductor subject to specified angular restrictions on their root numbers. Its Theorem 1.1 is not individual nonvanishing for every character. The distinction is essential: even a density-one result could leave exceptions. The source's introductory prose contains an apparent sign typo when describing an older nonvanishing proportion; no conclusion here is drawn from that sentence. The theorem's actual nonvanishing inequality and stated family were inspected.

\paragraph{2. Convergence on the positive real axis.}
For a nonprincipal character $\chi$ modulo $q$, the complete residue sum is zero. If $a$ is a unit with $\chi(a)\ne1$, multiplication by $a$ permutes the residues and gives $\sum\chi=\chi(a)\sum\chi$, proving the assertion. Consequently its partial sums are bounded by $q$. Dirichlet summation shows that
\[
 L(s,\chi)=\sum_{n=1}^{\infty}\chi(n)n^{-s}
\]
converges for every real $s>0$, and represents the analytically continued $L$-function there. More explicitly, if $B_N(x)=\sum_{N<n\leq x}\chi(n)$, then $|B_N(x)|\leq2q$, and summation by parts gives the safe remainder bound
\[
 \left|L(s,\chi)-\sum_{n\leq N}\chi(n)n^{-s}\right|
 \leq2qN^{-s}.
\]
One may improve the constant by starting at a complete block. The displayed bound is sufficient for a rigorous finite certificate when a computed partial sum is separated from zero by more than the error. It is not a uniform lower bound on the actual central value.

Regrouping into consecutive complete blocks preserves the sum. It does not license arbitrary rearrangements of the conditionally convergent series. The positivity arguments below use consecutive blocks only.

\paragraph{3. Complete elementary subcases.}
For the primitive quadratic character modulo three,
\[
 L(s,\chi_{-3})=\sum_{k\geq0}\bigl((3k+1)^{-s}-(3k+2)^{-s}\bigr)>0
 \qquad(s>0).
\]
Each block is positive. The same argument for the character modulo four gives
\[
 L(s,\chi_{-4})=\sum_{k\geq0}\bigl((4k+1)^{-s}-(4k+3)^{-s}\bigr)>0.
\]
The finite first block already makes the limit strictly positive; later positive blocks cannot cancel it.

There are three primitive characters modulo five. The real one has values $1,-1,-1,1$ on $1,2,3,4$. Set $f(x)=x^{-s}$. Since $f$ is strictly convex for $s>0$, its decrement $f(x)-f(x+1)$ strictly decreases with $x$. Therefore
\[
 f(5k+1)-f(5k+2)-f(5k+3)+f(5k+4)>0.
\]
This proves nonvanishing of the real character. For a complex character with $\chi(2)=i$, the values on $1,2,3,4$ are $1,i,-i,-1$. Hence
\[
 \operatorname{Re} L(s,\chi)
 =\sum_{k\geq0}\bigl((5k+1)^{-s}-(5k+4)^{-s}\bigr)>0.
\]
Its conjugate has the same positive real part. All primitive characters of conductor five are thus covered, including the nonreal ones.

For conductor eight, the two primitive characters have nonzero values on $1,3,5,7$ equal respectively to $1,-1,-1,1$ and $1,1,-1,-1$. Strict convexity gives positivity for the first pattern, by comparing decrements across intervals of length two. Termwise comparison gives positivity for the second. The remaining nonprincipal character modulo eight is induced from conductor four and need not be counted as primitive of conductor eight.

Finally, the conductor-one case is $\zeta(1/2)$. For $s>0$ the alternating eta series has the positive block representation
\[
 \eta(s)=\sum_{k\geq1}\bigl((2k-1)^{-s}-(2k)^{-s}\bigr)>0.
\]
The identity $\eta(s)=(1-2^{1-s})\zeta(s)$, initially proved for $s>1$ and continued to $s>0$, $s\ne1$, gives $\zeta(1/2)<0$. In particular it is nonzero. This handles a boundary case that is absent from nonprincipal-character arguments.

These proofs actually give nonvanishing along the entire positive real axis in the displayed nonprincipal cases. That stronger conclusion arises from their special residue patterns. General primitive real characters can have long, irregular sign patterns, and complex characters need not have block real parts of one sign.

\paragraph{4. What the functional equation does and does not force.}
Let $a\in\{0,1\}$ be determined by $\chi(-1)=(-1)^a$, and set
\[
 \Lambda(s,\chi)=\left(\frac q\pi\right)^{(s+a)/2}
 \Gamma\left(\frac{s+a}{2}\right)L(s,\chi).
\]
For primitive $\chi$ its functional equation is
\[
 \Lambda(s,\chi)=\varepsilon_\chi\Lambda(1-s,\overline\chi),
 \qquad |\varepsilon_\chi|=1.
\]
At the real central point, conjugation of the Dirichlet series and analytic continuation show that $w=\Lambda(1/2,\chi)$ satisfies $w=\varepsilon_\chi\overline w$. Write $\varepsilon_\chi=e^{i\theta}$. Then $e^{-i\theta/2}w$ is real. This confines $w$ to a line in the complex plane, but that line contains both zero and nonzero points. It proves neither nonvanishing nor vanishing.

In particular, when $\varepsilon_\chi=-1$ and $\chi$ is nonreal, the equation says $w$ is purely imaginary. The conclusion $w=-w$ would require the extra self-duality identity $\chi=\overline\chi$. For primitive real Dirichlet characters the classical quadratic Gauss-sum evaluation gives the completed root number $+1$ in this normalization. Thus the familiar forced central zero for a self-dual function with sign $-1$ does not automatically give a counterexample among these primitive Dirichlet functions. The Gauss-sum evaluation is a standard external input; the logical distinction between $w$ and $\overline w$ is proved directly above.

\paragraph{5. A general positivity criterion and its limitation.}
Suppose a real periodic sequence $c_n$ of mean zero has nonnegative initial sums $A_m=\sum_{n=1}^m c_n$ for all $m$, with at least one positive sum. For $s>0$, summation by parts gives
\[
 \sum_{n\geq1}c_n n^{-s}
 =\sum_{m\geq1}A_m\bigl(m^{-s}-(m+1)^{-s}\bigr)>0.
\]
The boundary term vanishes because $A_m$ is bounded. Every coefficient multiplying $A_m$ is positive. This is a useful sufficient condition that can be checked on one period. It is not necessary for nonvanishing and is not satisfied by all real characters. One can also apply it to the real part of a fixed rotation of a complex character, but finding a rotation that makes every partial sum nonnegative is an extra hypothesis, not a consequence of the root-number line condition.

For the conductor-five real character, some block arguments exploit convexity even when a simpler monotonicity criterion would be less informative. More generally, finite differences of $x^{-s}$ have alternating signs, suggesting structured groupings. The arithmetic gap is to obtain such a sign-controlled decomposition for every primitive character. Periodicity alone does not provide one: a general mean-zero periodic sequence can have its central Dirichlet sum zero through cancellation.

\paragraph{6. Why moments and computations leave an individual gap.}
For a finite family of central values $z_\chi$, Cauchy--Schwarz gives
\[
 \#\{\chi:z_\chi\ne0\}\geq
 \frac{|\sum_\chi z_\chi|^2}{\sum_\chi|z_\chi|^2},
\]
when the denominator is nonzero. Mollifiers alter the values but preserve the fact that a zero contributes zero. This can force many nonzero values, as in [S1], yet the inequality contains no assertion about a specified exceptional index. Replacing a positive proportion by every index is a missing theorem, not an improvement in notation.

A numerical value such as $10^{-20}$ is not a certificate without an error smaller than its distance from zero. A valid finite procedure can enclose each $n^{-1/2}$ by rational endpoints using integer square comparisons, sum these intervals with the character coefficients, and then add the proved tail bound. If the resulting interval or rectangle excludes zero, that individual value is certified. If it contains zero, the output is inconclusive, not a counterexample. No finite conductor cutoff settles the universal statement.

\paragraph{7. Diagnostics and outcome.}
The saved exact check verifies the character tables used here, multiplicativity on units, the conductor distinctions, and strict positivity certificates for the first grouped blocks using rational square-root intervals. The infinite positivity follows from the monotonicity and convexity proofs, not from checking finitely many blocks. The functional-equation diagnostic verifies the rotated-line identity on algebraic model values and does not claim to compute actual root numbers in general.

\textbf{UNRESOLVED.} All primitive characters of conductors $1,3,4,5,8$ are handled, with both complex characters at conductor five included. The first unrestricted gap is an individual lower bound or cancellation obstruction for arbitrary primitive characters. The restricted-family moment theorem and root-number symmetry do not supply that missing step.

\subsection{Sources}
\begin{itemize}
\item[S1]Non-vanishing of Dirichlet L-functions at the central point with restricted root number; arXiv2603.22124v1,23March2026.
\url{https://arxiv.org/pdf/2603.22124v1}.
\textit{Formulation source.}
\item[E1]A. Earnst, Non-Vanishing of Dirichlet L-functions at the central point with restricted root number, arXiv:2603.22124v1, Theorem 1.1 and Remark 1.2, HTML rendering.
\url{https://arxiv.org/html/2603.22124v1}.
\textit{Read 27 September 2026: positive-proportion result for even primitive prime-conductor characters in specified angular windows, not all-character nonvanishing.}
\end{itemize}
\end{document}
